%% Chapter 282 in the approved edition contents.  Editorial reference
%% navigation: registry files own added rule displays; source treatments and
%% their proof standings remain at the linked reader-facing destinations.
\chapter{Classical and Constructive Calculus Sheets}
\label{sec:p14-classical-and-constructive-calculi}
\chaptermark{Classical and Constructive Calculi}

Choose a presentation before checking a proof.  A logic's name does not fix
its formula grammar, derivation objects, treatment of assumptions, or equality
rules.  The sheets in this chapter retrieve the registered presentations;
the accompanying comparisons identify the inputs they accept and the
conditions that a displayed rule instance must satisfy.  The full
natural-deduction and sequent treatments are in
Chapters~\ref{ch:prop-deduction} and~\ref{ch:fol-deduction}; the
proof-theoretic comparison continues in Part~\ref{part:p07}
\cite{Gentzen1934a,Kleene1952,Prawitz1965}.

These lookup sheets do not replace the inherited rule displays, examples,
or proofs.  A link to a theorem retrieves its own statement and proof
standing: an inline source proof, a linked annex proof, a stated sketch, or
a cited result.  In particular, equivalence of consequences does not imply
identity of proofs, preservation of proof size, or preservation of normal
forms.  The presentations and translation hypotheses must be matched before
using an adequacy or normalization result.

\section{Propositional and first-order natural deduction}
\label{sec:p14-classical-calculi-natural-deduction}

The classical sheets use finite rooted formula trees, a single conclusion,
and primitive negation and falsity.  Introduction and elimination rules
govern the connectives; RAA supplies classical strength.  A supporting
assumption set records which formulas may occur at open leaves.  It may
contain unused formulas, but an unused permission is not an assumption
occurrence.  Each label discharged by the registered classical presentations
must mark at least one actual open occurrence in the relevant subderivation.
For disjunction elimination the two minor branches have separate discharge
classes; implication introduction, negation introduction and RAA have one.
There is no primitive repetition rule on either sheet.

When grafting a proof into several open leaves, use separately copied
derivations with labels fresh for the receiving derivation and for one
another.  Otherwise a rule above a graft can accidentally discharge an
assumption in the inserted proof.  The exact propositional operation and
its proof are Lemma~\ref{lem:nd-substitution}.  This replacement of an
assumption by a proof is different from substituting an individual term
into a first-order formula.

\paragraph{Propositional lookup.}
Use Sheet~\ref{sheet:nd-cpl} for the full primitive propositional language.
The rule applications and worked excluded-middle proof are in
Section~\ref{sec:prop-deduction-01}.  Propositional soundness is
Theorem~\ref{thm:prop-nd-soundness}; the comparison with Hilbert and sequent
derivability is Theorem~\ref{thm:p02-calculi-derivability-equivalence}.
The latter is a cited equivalence and explicitly translates the full
language into the Hilbert fragment; it is not a step-by-step conversion
claim.

\showcalculus{nd-cpl}

\paragraph{First-order lookup.}
Sheet~\ref{sheet:nd-fol} retains the entire classical propositional surface
and adds quantifier rules, explicit bound-variable alpha-conversion, and
logical identity.  The syntax uses the total capture-avoiding operation of
Definition~\ref{def:p02-fol-capture-avoiding-substitution}; bound names are
renamed when required.  Thus universal elimination and existential
introduction do not carry a separate partial-substitution definedness
test.  Universal introduction nevertheless requires that the generalized
variable be free in no open assumption.  Existential elimination requires
that its witness variable be absent from the conclusion, every other open
assumption, and the existential major premise.  Capture avoidance does not
remove either eigenvariable condition.

The source's universal-reasoning example uses a fresh variable and then
the explicit alpha rule to reach the requested quantified string.
The existential-elimination example keeps its witness inside the minor
derivation; both are in Section~\ref{sec:fol-deduction-01}.
Alpha-conversion changes bound names only.  It is neither equality
elimination nor permission to replace a formula by any classically
equivalent formula.  Identity introduction and replacement have their
separate rules; their derived symmetry and transitivity are retrieved in
Section~\ref{sec:p14-classical-calculi-equality} below.

\showcalculus{nd-fol}

For this first-order natural deduction, soundness and its complete inline
proof are Theorem~\ref{thm:fol-soundness}.  Strong completeness is
Theorem~\ref{thm:fol-completeness}, with the arbitrary-signature construction
and its choice assumptions stated there.  Neither result supplies a
normalization theorem for classical RAA or for an unnamed equality
extension of intuitionistic natural deduction.

\section{Sequent, Hilbert, tableau, and resolution systems}
\label{sec:p14-classical-calculi-proof-systems}

These four interfaces encode different questions.  A sequent derivation
establishes a left-to-right consequence with explicit structural data.
A Hilbert derivation ends a finite formula sequence.  A closed tableau
refutes its root formulas, and a resolution derivation refutes clauses by
reaching the empty clause.  To prove a classical consequence by refutation,
include the negated proposed conclusion in the root or clausified input.
The conversion to clauses must be recorded separately: Boolean equivalence
and Skolem equisatisfiability are not the same guarantee
\cite{Smullyan1968,Robinson1965,BachmairGanzinger2001}.

\subsection{Two classical sequent presentations}

Sheet~\ref{sheet:lk} has primitive negation, contexted identity initials,
and finite-multiset contexts on both sides.  Commas add multiplicities.
Exchange is built into those multisets; weakening and contraction remain
explicit rules.  The succedent is disjunctive, including the empty
succedent, and is not a single-conclusion natural-deduction context.
Its propositional restriction omits the four quantifier rows.  Its
first-order equality-free extension keeps them, with eigenvariables free
in no formula of the conclusion sequent for universal right and
existential left.  The original displays are in
Sections~\ref{sec:prop-deduction-02} and~\ref{sec:fol-deduction-02}.

\showcalculus{lk}

Sheet~\ref{sheet:lk-defneg} is a different presentation: negation
abbreviates implication to falsity, identity is uncontexted, and
implication left combines separate premise contexts.  It has no primitive
negation or equality rows.  The canonical-to-defined-negation map is
Definition~\ref{def:p07-defined-negation-translation}; its derivability
equivalence is Theorem~\ref{thm:p07-lk-presentation-equivalence}, whose
annex proves equivalence of derivability without identifying proof objects.

\showcalculus{lk-defneg}

Canonical LK soundness is proved inline in
Theorems~\ref{thm:prop-lk-soundness} and~\ref{thm:fol-lk-soundness}.
The separate complete finite-sequent proof is
Theorem~\ref{thm:p07-canonical-lk-completeness}, with
Annex~\ref{annex:p7.calculi.4}; its first-order form is equality-free and
distinguishes arbitrary from countable signatures.  Its ND bridge is
Theorem~\ref{thm:p07-lk-nd-equivalence}, with
Annex~\ref{annex:p7.calculi.3}.  That bridge translates the succedent to a
disjunction and the antecedent to its support, including an empty
succedent, and uses Cut.  It does not erase these steps merely because the
resulting consequences agree.

For the exact defined-negation sheet, the completed Cut-and-Mix
elimination construction is Theorem~\ref{thm:p07-classical-cut-elimination}
and Annex~\ref{annex:p07.classical.cut-elimination}.  It preserves the
literal endsequent and context multiplicities.  That theorem is not a
cut-elimination theorem for primitive-negation LK or for
Sheet~\ref{sheet:lk-eq}.

\subsection{Hilbert sequences and dependency restrictions}

Sheet~\ref{sheet:hilbert-cpl} uses the primitive
\(\neg,\to\) fragment, axiom-schema instances, assumptions, and modus
ponens.  Uniform substitution instantiates a schema; it is not a rule
permitting arbitrary alteration of a premise line.  Use
Definition~\ref{def:p02-hilbert-translation} before comparing this fragment
with the other full-language propositional sheets.  The deduction theorem
has its complete proof at Theorem~\ref{thm:p02-hilbert-deduction} and
Annex~\ref{annex:p2.calc.2} \cite{Mendelson1997}.

\showcalculus{hilbert-cpl}

Sheet~\ref{sheet:hilbert-fol} is equality-free and retains primitive
existential quantification.  Its registered editorial presentation combines
Mendelson's universal Hilbert base with Kleene's existential-introduction
schema and premise-sensitive existential-elimination rule; it is not
attributed as an unchanged calculus of just one of those sources
\cite{Mendelson1997,Kleene1952,Enderton2001}.
Definition~\ref{def:p02-hilbert-fol-dependencies} records a finite dependency
set on every line.  Generalization is licensed only when its variable is
free in none of the premises on which that line depends.  The existential
elimination rule also excludes the variable from its consequent.  The
universal distribution schema separately excludes its bound variable from
the antecedent formula.  These are three different checks, not an
unrestricted instruction to quantify a previously written line.

\showcalculus{hilbert-fol}

The premise-sensitive deduction theorem is
Theorem~\ref{thm:p02-hilbert-fol-deduction}, with
Annex~\ref{annex:p2.fol.calc.5}. Full adequacy of the literal registered
presentation is supplied separately at
\cref{sec:p02-hilbert-fol-adequacy,thm:p02-hfa-adequacy}.
The proof preserves dependency-sensitive Gen and existential elimination,
derives alpha-conversion, and uses a closed contradiction for its
equality-free full-language bridge. Equality is excluded throughout the
translated proof tree. Finite-input, countable-signature and arbitrary
signature cases retain their distinct set-theoretic assumptions; an
atomless signature requires the explicitly declared predicate expansion.

\subsection{Retained-formula analytic tableaux}

Both tableau sheets retain principal formula occurrences after expansion.
A propositional occurrence is marked, so a descendant branch cannot expand
that same occurrence again.  A split copies the retained branch and its
marks to both children.  A branch closes on falsity or a complementary
formula pair; closing one branch does not close the whole tableau.
The propositional adequacy statement and saturated-open-branch valuation
are Theorem~\ref{thm:p02-tableau-adequacy}.  Its soundness, termination,
truth-lemma and completeness clauses are proved in
\cref{sec:p02-w02-p02-014-tableau-metatheory}, which also shows that
propositional tableaux decide satisfiability
(\cref{thm:p02-w02-p02-014-decision}).

\showcalculus{tableaux-cpl}

In the first-order sheet, gamma occurrences remain available at every
unused closed term of the specified parameter-extended root language.
Delta occurrences expand once using the least parameter not previously
on that branch.  Thus the two quantifier expansion kinds do not share
the one-use propositional policy.  Fairness ranges over the entire fixed
closed-term enumeration, not only terms already visible on a branch.
The exact finite equality-free sentence input, finite closed-refutation
equivalence, and fair saturated open-branch model construction are
Definition~\ref{def:p02-fol-tableau} and
Theorem~\ref{thm:p02-fol-tableau-adequacy}, proved in
\cref{sec:p02-w02-p02-014-fol-herbrand,sec:p02-w02-p02-014-fol-systematic},
where equality rules are also added; an unfinished or unfair search
supplies no countermodel (\cref{def:p02-w02-p02-014-fol-unfair}).

\showcalculus{tableaux-fol}

\subsection{Clause refutation}

Propositional clauses are finite sets of literals, not lists of formulas
whose multiplicities survive an inference.  Resolution removes the selected
complementary literals and unions the remainders.  The resolvent may be
tautological or empty.  The inherited soundness proof is inline in
Section~\ref{sec:prop-metatheory-04}; finite refutational completeness is
Theorem~\ref{thm:prop-resolution-complete}, and the arbitrary-set extension
has its separate finiteness statement at
Remark~\ref{rem:resolution-finiteness}.

\showcalculus{resolution-cpl}

First-order clauses denote universal closures.  Before every binary
resolution inference, take fresh variable-renamed copies of the two parent
occurrences with disjoint variable sets, even when both parents come from
one stored clause.  Resolve atoms with a most general unifier, apply the
substitution simultaneously to the entire remainder, and retain Factoring
for same-polarity literals.  Identified duplicates disappear by set equality.
The exact rule displays and full soundness proof are in
Section~\ref{sec:fol-cardinality-normalforms-04}.  The inherited
Theorem~\ref{thm:fol-resolution} is a finite equality-free completeness
sketch, using the cited Herbrand and lifting results
(Theorems~\ref{thm:herbrand-clauses} and~\ref{thm:resolution-lifting}); a
fresh constant is added when the signature has no ground term.  The lookup
does not promote those dependencies to locally reproduced proofs.

\showcalculus{resolution-fol}

\section{Equality, substitution, and unification rules}
\label{sec:p14-classical-calculi-equality}

Three operations must be kept apart.  Formula substitution replaces free
individual variables with terms while avoiding capture.  Equality
replacement uses a proved identity to transport a formula instance.
Unification solves a metalinguistic syntactic equation problem by finding
a substitution.  The fact that all three involve terms does not make them
interchangeable.

\begingroup\small
\begin{longtable}{@{}p{.20\textwidth}p{.33\textwidth}p{.39\textwidth}@{}}
\caption{Substitution and equality interfaces.}
\label{tab:p14-classical-calculi-substitution}\\
\toprule
Operation & Required data & Controlling destination and limit \\
\midrule
\endfirsthead
\caption[]{Substitution and equality interfaces (continued)}\\
\toprule
Operation & Required data & Controlling destination and limit \\
\midrule
\endhead
\midrule\multicolumn{3}{r}{\emph{continued on next page}}\\
\endfoot
\endlastfoot
Formula substitution & A formula, individual variable, and term; rename
bound variables to prevent capture. &
Definition~\ref{def:p02-fol-capture-avoiding-substitution}.  This is a total
syntax operation, not an inference from a formula to an arbitrary instance. \\
Assumption substitution & A selected family of open assumption occurrences
and a derivation of their formula; separately freshen inserted copies. &
Lemma~\ref{lem:nd-substitution}; for quantified empty-discharge NJ,
Lemma~\ref{thm:p07-nj-substitution} and Annex~\ref{annex:p7.nj.substitution}.
The latter includes internal eigenvariable renaming. \\
Clause substitution & A finite-support simultaneous variable-to-term map;
apply it once to the original literals. &
Section~\ref{sec:fol-cardinality-normalforms-04} fixes action and composition.
It does not recursively substitute inside an image inserted in that same
application. \\
Logical equality replacement & An identity premise and a matching formula
instance, using capture-avoiding substitutions. &
Sheet~\ref{sheet:nd-fol}, Theorem~\ref{thm:p02-equality-laws}, and
Annex~\ref{annex:p2.fol.calc.4}.  Identity is semantic identity, not an
arbitrary reflexive relation. \\
Syntactic unification & Finite first-order terms in a free ranked term
signature; an equation multiset. &
Sheet~\ref{sheet:unification-first-order} and
Theorem~\ref{thm:p12-unification-mgu}.  No associative, commutative,
higher-order, or infinite-term theory is implicit. \\
\bottomrule
\end{longtable}
\endgroup

\paragraph{Equality-bearing sequents.}
Sheet~\ref{sheet:lk-eq} explicitly extends canonical LK, not the
defined-negation sheet, by identity reflexivity and left replacement.
The opposite replacement direction is derived using symmetry and Cut.
The full source derivation is in Section~\ref{sec:fol-deduction-02}; it
does not establish that the derived instance is cut-free.  Its soundness
proof is Theorem~\ref{thm:fol-lk-soundness}. The exact cut-admitting
ND equivalence and completeness are now
\cref{thm:p07-eqlk-nd,thm:p07-eqlk-completeness}. For this literal
surface, Cut and Mix are not admissible in the fragment omitting both:
\cref{thm:p07-eqlk-cut-failure} gives the essential-Cut symmetry
derivation and a complete backward-branch obstruction. Pure-LK or
defined-negation cut elimination must not be transferred to it.

\showcalculus{lk-eq}

\paragraph{Finite-term unification.}
The inherited accumulator procedure and its proof sketch are in
Section~\ref{sec:fol-cardinality-normalforms-04}.  The separately registered
Martelli--Montanari transformation interface keeps Delete, Decompose,
Clash, Orient, Occurs and Eliminate distinct.  A successful solved form
determines a most general substitution, in the factorization sense of
Theorem~\ref{thm:p12-unification-mgu}, with source attribution
\cite{MartelliMontanari1982,BaaderSnyder2001}. The complete local literal-sheet
proofs are \cref{thm:p12-ue-sheet-invariant,thm:p12-ue-sheet-termination,
thm:p12-ue-extraction}; the separate accumulator has the explicit composition
invariant of \cref{thm:p12-ue-accumulator}.
The source accumulator and the registered solved-equation state are not
identical algorithms merely because they solve the same finite syntactic
problem.

The sheet retains an eliminated equation and extracts all final bindings
simultaneously. The inherited accumulator removes the chosen binding and
updates its substitution. Their intermediate states need not agree.
\Cref{thm:p12-ue-renaming} concerns most general solutions up to finite
variable renaming, not uniqueness of factor substitutions; the latter
can fail even for the one-equation problem \(x\doteq a\).

The occurs check rules out a binding such as \(x\mapsto f(x)\) as a
solution of \(x\doteq f(x)\) by finite terms; it does not prohibit
writing that finite-support substitution as syntax.  See
Example~\ref{def:p12-occurs-check}.  Dropping the check changes the term
domain and no longer uses the stated finite-term theorem.

\showcalculus{unification-first-order}

\paragraph{Equality in automated refutation.}
Ordinary resolution permits equality atoms for soundness but has no
identity-complete inference family.  Ordered superposition instead requires
a specified simplification ordering, negative-literal selection,
post-substitution eligibility, sound redundancy and fair saturation.
The three registered inference families are Superposition, Equality
Resolution and Equality Factoring.  Their ordering and eligibility
conditions must be checked on the instantiated literals; replacement is
not unrestricted rewriting at any chosen occurrence.
Definition~\ref{def:p12-ordered-superposition} specifies the saturation
interface, and Theorem~\ref{thm:p12-superposition-completeness} is the
cited completeness boundary \cite{BachmairGanzinger2001,NieuwenhuisRubio2001}.
It cannot be transferred to a search that drops one of those hypotheses.

\showcalculus{ordered-superposition}

\section{Intuitionistic, minimal, intermediate, and typed calculi}
\label{sec:p14-classical-calculi-constructive}

\subsection{Intuitionistic sequents and assumption discharge}

Sheet~\ref{sheet:lj} has exactly one succedent formula, no right weakening
or right contraction, defined negation, and no equality.  These are
restrictions of the derivation language, not cosmetic omissions in a
two-sided classical display.  Its propositional fragment is the calculus
named on the IPL card, Entry~\ref{entry:p08-ipl}.  The complete local
Kripke adequacy proof is in Section~\ref{sec:p08-ipl-kripke-completeness}.
For first-order expanding-domain semantics, use
Definition~\ref{def:p08-iqc-expanding-domain} and the cited theoremhood
result, Theorem~\ref{thm:p08-iqc-completeness}.  The latter does not state
an arbitrary-premise first-order completeness theorem.

The exact Cut-and-Mix elimination theorem for LJ is
Theorem~\ref{thm:p07-lj-exact-cut-elimination}, with
Annex~\ref{annex:p7.lj.exact-cut-elimination}.  It includes contraction,
the sheet's split-context implication left, and the quantifier conditions;
single-succedent syntax remains intact throughout the transformation.

\showcalculus{lj}

The two NJ sheets have the same equality-free defined-negation formula
grammar and eigenvariable conditions.  Sheet~\ref{sheet:nj} requires
nonempty discharged assumption classes.
Sheet~\ref{sheet:p07-nj-vacuous} permits an empty class in implication
introduction, each minor disjunction branch, and the minor existential
derivation.  Those permissions must not be inferred merely by writing an
unused assumption in a supporting set.

The full first-order normalization theorem is for the empty-discharge
presentation: Theorem~\ref{thm:p07-nj-vacuous-normalization} and
Annex~\ref{annex:p7.nj.vacuous-normalization}.  It supplies a terminating
declared strategy with the same conclusion and no new open assumption
formula; it is not a claim that every interleaving of all conversion
classes terminates.  The consequence bridge and nonvacuous obstruction
are separately proved in
Theorems~\ref{thm:p07-nj-vacuous-consequence-bridge}
and~\ref{thm:p07-nj-nonvacuous-obstruction}, with their reciprocal annex
proofs.  Translating a normalized proof back can insert detours, so equal
consequences do not supply a normal form in the nonvacuous calculus
\cite{Prawitz1965,DeGroote2002NJPermutation,Read2014ProofTheoreticValidityManuscript}.

\showcalculus{nj}
\showcalculus{p07-nj-vacuous}

\subsection{Minimal and intermediate boundaries}

For a minimal natural-deduction restriction, keep the stated
defined-negation NJ grammar and all of its chosen discharge and
eigenvariable conditions, but omit falsity elimination.  Choosing
nonvacuous or empty discharge still matters: ``minimal'' does not settle
that presentation choice.  Removing RAA alone from classical ND retains
falsity elimination and therefore does not make the result minimal.
Moreover, the primitive-negation classical sheet is not literally the
defined-negation NJ syntax.

Sheet~\ref{sheet:minimal-nd} now fixes the empty-discharge propositional
minimal presentation. Its arbitrary persistent-falsum semantics and full
strong local adequacy are \cref{thm:p08-mi-minimal-adequacy}; the separate
positive-language/falsum-recovery bridge is
\cref{thm:p08-mi-minimal-falsum-recovery}. A step from falsity to an
arbitrary formula uses the omitted rule. The Kuroda embedding in
\cref{thm:p08-kuroda-embedding} still names LJ, not this minimal target.
Johansson's primitive-negation presentation is a different source obligation.

\showcalculus{minimal-nd}
\showcalculus{hilbert-kc}
\showcalculus{hilbert-lc}

Intermediate propositional logics are theorem sets between IPL and
classical propositional logic, closed under uniform substitution and
modus ponens, as fixed in Definition~\ref{def:p08-intermediate-logic}.
This definition does not choose a Hilbert, sequent, natural-deduction or
hypersequent rule surface.  The two current concrete profiles are
KC, IPL plus weak excluded middle \(\neg A\lor\neg\neg A\), and LC,
IPL plus prelinearity \((A\to B)\lor(B\to A)\).
Theorems~\ref{thm:p08-kc-profile} and~\ref{thm:p08-lc-profile} retain
their cited sources, \cite{Jankov1968WeakExcludedMiddle,Dummett1959}.
The matched modern sheets above have complete strong frame adequacy,
finite countermodel decisions and algebraic/chain bridges in
\cref{thm:p08-mi-intermediate-adequacy,thm:p08-mi-finite-decision,
thm:p08-mi-kc-algebra-bridge,thm:p08-mi-finite-chain-bridge}; these do not
replace the separately attributed original matrix/algebra presentations.
Table~\ref{tab:p08-intermediate-logics} distinguishes their frame classes,
and the adjoining countermodels separate their strict theorem-set
inclusions.  Both fail the theorem-level disjunction property, although
the IPL result is retained in
Theorem~\ref{thm:p07-ipc-disjunction-property}.  No dedicated KC or LC
rule sheet is silently manufactured by adjoining these formulas to an
unspecified base calculus.

\subsection{Types, explicit polymorphism, and dependency}

The arrow-only STLC sheet has finite contexts assigning distinct term
variables to simple types.  Abstraction extends the context by a fresh
term variable; that binder may have no occurrence in the body.  Its type
arrow is not an object-language formula implication, and the sheet has no
products, sums, polymorphism or dependent types.  The larger
Curry--Howard categorical extension is defined separately in
Definition~\ref{def:p07-curry-howard}; its extra constructors and
beta--eta equalities do not alter this sheet.  Complete typing substitution,
subject reduction and strong normalization are
Theorems~\ref{thm:p07-stlc-substitution},
\ref{thm:p07-stlc-subject-reduction} and
\ref{thm:p07-stlc-strong-normalization}, with their annex proofs.

\showcalculus{stlc}

For comparison, untyped lambda calculus has no native typing judgment.
Compatible beta reduction may act under a binder, and capture avoidance
remains compulsory.  The Church--Rosser theorem and complete annex proof
(Theorem~\ref{thm:p07-church-rosser}, Annex~\ref{annex:p7.lambda.1})
give confluence, not termination: the self-application loop in
Definition~\ref{def:p07-beta-conversion} has no normal form.
The S--K sheet instead has application and two primitive contractions but
no binder.  The bracket-abstraction construction and proved operational
bridge are Definition~\ref{def:p07-sk-combinators} and
Theorem~\ref{thm:p07-bracket-abstraction}; they translate syntax rather
than identifying lambda terms with combinatory terms
\cite{ChurchRosser1936,CurryFeys1958,Barendregt1984}.

\showcalculus{lambda-untyped}
\showcalculus{sk-combinatory}

Church-style System F adds a separate type-variable context, explicit
type abstraction and application, and universal type formation.
Universal type introduction excludes its type variable from every type
assigned in the term context, as well as requiring freshness for the
type-variable context.  Term substitution and type substitution have
different domains; both avoid capture.  The full local confluence,
strong-normalization and explicit-checking results are
Theorems~\ref{thm:p07-system-f-confluence},
\ref{thm:p07-system-f-strong-normalization} and
\ref{thm:p07-system-f-church-checking}.
Erasing the annotations produces the separate Curry-style assignment
problem, whose cited undecidability result is
Theorem~\ref{thm:p07-system-f-curry-undecidable}; the explicit checking
algorithm is not an algorithm for that erased input
\cite{Girard1972,GirardLafontTaylor2003Web,Wells1999}.

\showcalculus{system-f}

System T extends the arrow-only typed terms by zero, successor and
recursors at finite types. The argument order of each recursor is binding
on the sheet. Distinguish its term calculus from the quantifier-free
equality theory in \cref{def:p07-dialectica-target}. Extracting a term
and proving its matrix property are the two conclusions of
\cref{thm:p07-dialectica-soundness}.
\Cref{thm:p07-system-t-numerical-computability} supplies the separate
proof that closed numerical terms can be evaluated effectively. It
handles the recursors explicitly; STLC's inert-constant corollary alone
does not cover their computation rules \cite{AvigadFeferman1998}.

\showcalculus{system-t}

The selected intensional MLTT sheet distinguishes context, type, term,
type-equality and term-equality judgments.  A substitution through a
dependent context must also substitute into its suffix and judgment.
Judgmental computation and conversion are different from inhabiting an
identity type: the sheet adds no equality reflection.  Its dependent
products, sums, identity types and cumulative Russell universes constitute
an explicitly identified editorial composite
\cite{MartinLof1984,NordstromPeterssonSmith1990,Hofmann1997}.
The cited Abel--Coquand--Dybjer normalization and typed-equality theorem,
Theorem~\ref{thm:p07-mltt-normalization-decision}, controls only dependent
functions and one Russell universe, not the composite's sums, identity
types or cumulative hierarchy.  The CoC theorem and CIC boundary are
likewise separate at Theorem~\ref{thm:p07-coc-strong-normalization} and
Boundary~\ref{def:p07-cic-pointer}.  No blanket normalization or checking
claim for all dependent type theories is obtained by this lookup.

\showcalculus{mltt-intensional}

\section{Source variants and side-condition comparisons}
\label{sec:p14-classical-calculi-variants}

The controlling unit is a specified presentation, not the shortest
familiar system name.  The inherited displays keep their source grammar;
registered lookup sheets make their own additional policies explicit.
The Hilbert first-order and MLTT composites are identified above rather
than attributed as unchanged single-source systems.  A formula translation
or consequence bridge licenses only its stated transfer.  It cannot
silently transfer equality, a structural rule, normalization, or a
proof-search policy.

\begingroup\small
\begin{longtable}{@{}p{.20\textwidth}p{.34\textwidth}p{.38\textwidth}@{}}
\caption{Conditions that distinguish neighboring calculus presentations.}
\label{tab:p14-classical-calculi-side-conditions}\\
\toprule
Comparison & Condition to inspect & Consequence of changing it \\
\midrule
\endfirsthead
\caption[]{Conditions that distinguish neighboring calculus presentations (continued)}\\
\toprule
Comparison & Condition to inspect & Consequence of changing it \\
\midrule
\endhead
\midrule\multicolumn{3}{r}{\emph{continued on next page}}\\
\endfoot
\endlastfoot
ND assumptions / sequent contexts & Actual labeled open leaves and
supporting sets versus formula multiplicities in multisets. & An unused
context permission cannot satisfy nonvacuous discharge.  Forgetting
multiplicity in an ND translation does not preserve a sequent proof tree. \\
Canonical LK / defined-negation LK & Primitive negation and contexted Id
versus defined negation, uncontexted Id, and split implication left. &
Use the explicit translation in
Definition~\ref{def:p07-defined-negation-translation}; the completed
defined-negation Cut theorem does not cover the primitive-negation sheet. \\
LK / LJ & Multiple or empty succedent versus exactly one succedent;
right structural rules versus none. & Classical right rules and
multi-succedent proof steps cannot be copied into LJ.  Classical and Kripke
adequacy statements use different semantic clauses. \\
Classical ND / quantified NJ & Primitive negation, RAA and identity versus
defined negation, no RAA, and no identity rules. & Omitting one classical
rule does not literally translate the grammar or prove normalization
for an equality extension. \\
Nonvacuous / empty-discharge NJ & Whether every carried label marks an
actual open assumption occurrence. & The consequence bridge preserves
conclusions and assumptions, not normality.  Branch-deleting conversions
can make a discharge class empty. \\
Formula substitution / eigenvariables & Capture-avoiding insertion versus
freshness from specified contexts and conclusions. & Renaming bound
variables to avoid capture does not license generalization from an open
\(P(x)\) assumption or export an existential witness. \\
Propositional / first-order tableaux & One-use marks versus repeated gamma
instances and fresh one-use delta parameters. & Expanding each universal
only once, or checking only already visible terms, is not fair saturation
for the first-order theorem. \\
Resolution / identity refutation & Renaming apart, MGU and Factoring versus
the additional ordering, equality-inference and saturation policies. &
An equality-bearing empty-clause derivation can be sound without an
equality-complete search procedure. \\
Church / Curry System F & Explicit binder types, type abstraction and
type application versus erased term syntax. & The input and decision
problem change; explicit checking and erased typability have different
decidability standings. \\
MLTT judgmental / identity equality & Conversion by judgmental equality
versus a term inhabiting an identity type. & No equality-reflection step is
licensed; a theorem for the narrower one-universe function core does not
certify the whole composite sheet. \\
\bottomrule
\end{longtable}
\endgroup

For a proposed rule instance, first match its formula or type grammar and
context representation; then check every premise, the actual discharge
classes, and the local freshness or eligibility condition.  Finally follow
the theorem link and match its semantics, input extent, equality standing,
and proof status.  This order prevents a correct theorem for a neighboring
presentation from being used as a certificate for the wrong proof object.
Full-source treatments, the separate proof constructions, and boundaries
for presentations not supplied by a sheet remain necessary after the
lookup succeeds.
